%---------------------------------------------------------------------
\chapter{Requirements and System Analysis}
%---------------------------------------------------------------------

\section{Stakeholders and Users}
The primary user is a \emph{compliance officer} at a producer, who assembles evidence and prepares
the working paper behind a return. A \emph{reviewer}, typically a senior officer or an internal
auditor, examines the working paper and accepts, qualifies or rejects each finding. The producer's
management relies on the reviewed result. The regulator (CPCB and the State Pollution Control
Boards) is affected indirectly, through the quality of what is filed, but is not a user: the
system does not interact with the portal.

\section{Functional Requirements}
\begin{description}[itemsep=1pt,font=\normalfont\bfseries]
\item[FR1] The system shall accept original files (PDF, CSV, JSON), hash their bytes, store them
immutably and issue a receipt for every intake outcome, including rejections.
\item[FR2] The system shall extract observations with exact locators (page, row, field) and record
every extraction failure.
\item[FR3] The system shall normalise units, metals and entity aliases with validated decimal
quantities, keeping missing, unknown, not applicable and zero distinct.
\item[FR4] The system shall freeze a case snapshot and a corpus snapshot, each identified by a
digest, before any investigation runs.
\item[FR5] The system shall reconcile obligation buckets against supported allocations with
deterministic tools, preserving credit conservation and earlier-liability priority.
\item[FR6] The system shall run a bounded investigation in which a model proposes one audited
action at a time, assesses the observation, and resolves, continues or abstains.
\item[FR7] The system shall accept a model-proposed finding only through a deterministic gate that
checks it against computed evidence, and shall record every rejected proposal with its reason.
\item[FR8] The system shall produce a working paper whose figures and citations are verified
against the snapshot, and shall bind each human decision to the exact report revision.
\item[FR9] The system shall export a run with its digests so that it can be replayed and checked.
\end{description}

\section{Non-Functional Requirements}
\begin{description}[itemsep=1pt,font=\normalfont\bfseries]
\item[NFR1] Exactness: compliance quantities use decimal arithmetic; binary floating point and
SQLite \texttt{REAL}/\texttt{SUM} are never used for balances.
\item[NFR2] Bounded cost: at most 12 model calls, 600~s of wall time and 32{,}768 tokens per model
call in one investigation; exhausted budgets are recorded as outcomes.
\item[NFR3] Traceability: every model call and tool action is recorded in an append-only audit
log before dispatch, with per-attempt records.
\item[NFR4] Safety: the model cannot approve evidence, impersonate a reviewer, modify originals,
run arbitrary code or SQL, or invent a source; dossier text is treated as untrusted data.
\item[NFR5] Isolation: hidden answers, fault labels and generator seeds never enter the runtime,
prompts, retrieval index, user interface or run artifacts.
\item[NFR6] Privacy: inference is local, on a loopback interface; no dossier text leaves the
machine during an investigation, and tokens and passwords are kept out of logs.
\item[NFR7] Portability: the base application and its tests run on Windows, Linux and macOS;
local inference dependencies are kept separate.
\end{description}

\section{Data Sources}
Table~\ref{tab:data} lists the data used. The two regulatory documents are public CPCB
publications kept with their SHA-256 digests; their redistribution rights are recorded as
\texttt{REVIEW\_REQUIRED}, so the public release lists them by title, digest and page count rather
than redistributing them. All dossiers are synthetic and watermarked ``SYNTHETIC -- NOT A CPCB
RECORD''.

\begin{table}[htbp]
\centering
\caption{Data sources used in ReCycleGuard}
\label{tab:data}
\small
\begin{tabularx}{\linewidth}{@{}l X l l@{}}
\toprule
\textbf{Source} & \textbf{Description} & \textbf{Size} & \textbf{Origin} \\ \midrule
R02 & E-Waste (Management) Rules, 2022~\cite{rules2022} (PDF) & 38 pages & CPCB website \\
R08 & CPCB FAQ under the E-Waste Rules, 23 Jan 2024~\cite{faq2024} (PDF) & 13 pages & CPCB website \\
\texttt{pilot-0.1} & Dossiers of PDFs, CSVs and JSON with hidden answers & 48 dossiers & Synthetic, team generator \\
Earlier sets & Hand-built and S06 development dossiers (development only) & -- & Synthetic, team-built \\
\bottomrule
\end{tabularx}
\end{table}

The \texttt{pilot-0.1} set is split into 24 development dossiers (a tuning set of 16 and a
checking set of 8) and 24 held-out dossiers. The held-out dossiers are not in the repository; a
seal file stores one SHA-256 per dossier so that their bytes can be checked without opening them.
Chapter~7 describes the set in detail.

\section{Hardware and Software Requirements}
Development used an Apple M3 laptop (where the early native experiments ran with MLX), a Windows
PC with an NVIDIA RTX GPU with 24~GB of memory (MG1, one Docker container per trial), and free
Kaggle notebooks with an NVIDIA Tesla T4 GPU with 16~GB (MG2, while the PC was unavailable). The
software is Python~3.11 or later managed with \texttt{uv}, LangGraph for the workflow, FastAPI for
the HTTP interface, a React user interface, and llama.cpp build \texttt{b12033} serving
Qwen3.5-9B quantised to Q8\_0~\cite{qwen35gguf,llamacpp}. Versions are listed in
Table~\ref{tab:stack} and Appendix~\ref{app:run}.

\section{Feasibility}
Inference is local, so there is no per-call charge; the cost is GPU time. A single-agent trial
used about 6 model calls and 90 to 110~s on a T4 with three parallel slots (Table~\ref{tab:mg2-rules-on}),
and the 448-trial rules-on queue of seven sessions took about three hours (log entry 14). Kaggle's free tier gives about 30 GPU
hours a week, which bounded the T4 series to three queues and one held-out run. Real producer data
and portal access were out of reach, which is why the benchmark is synthetic and the system never
connects to the portal.
